\subsubsection*{SUN}

The SUN (Showa University and Nagoya University) dataset is one of the largest publicly available colonoscopy video datasets, specifically designed to benchmark real-time AI-assisted polyp detection systems. It comprises 152{,}560 video frames, of which 49{,}799 contain visible polyps and 102{,}761 are negative samples. The data were extracted from 1405 colonoscopy procedures, with 100 independent polyps included in the final selection.

All frames were annotated with bounding boxes by three research assistants, then validated and corrected by two expert endoscopists. Additionally, a quality assurance process involving a blinded review by an external committee of experienced gastroenterologists was implemented, supervised by a contract research organization to ensure labeling integrity.

The dataset includes metadata for each polyp, including size, shape (e.g., protruded or flat), and pathology (e.g., low-grade or high-grade adenoma, carcinoma). Videos were recorded using Olympus high-definition endoscopes, ensuring consistency in image quality.

The SUN dataset was used to develop and evaluate an AI system based on the YOLOv3 architecture, optimized for real-time detection ($\approx$30 FPS). The system achieved 98.0\% sensitivity at the polyp level and 90.5\% sensitivity with 93.7\% specificity at the frame level. The authors emphasize the need for realistic and standardized benchmarks and explicitly designed SUN to fulfill this role, making it a critical resource for assessing detection performance under clinical constraints.
